% Compile with pdflatex twice, or run tectonic group_wl.tex.
\documentclass[11pt]{article}
\usepackage[T1]{fontenc}
\usepackage[utf8]{inputenc}
\usepackage{lmodern}
\usepackage[margin=1in]{geometry}
\usepackage{amsmath,amssymb,amsthm,mathtools}
\usepackage{microtype}
\usepackage{enumitem,booktabs}
\usepackage{xcolor}
\definecolor{linkink}{HTML}{193F60}
\usepackage[colorlinks=true,linkcolor=linkink,citecolor=linkink,urlcolor=linkink]{hyperref}
\hypersetup{pdftitle={EXPTIME-completeness of group Weisfeiler--Leman identification and equivalence},pdfsubject={EXPTIME-completeness with an explicit dense reduction}}
\newcommand{\F}{\mathbb F_p}
\newcommand{\eps}{\varepsilon}
\DeclareMathOperator{\Hom}{Hom}
\DeclareMathOperator{\Alt}{Alt}
\DeclareMathOperator{\im}{im}
\DeclareMathOperator{\rank}{rank}
\DeclareMathOperator{\rad}{rad}
\DeclareMathOperator{\spanop}{span}
\newtheorem{theorem}{Theorem}[section]
\newtheorem{lemma}[theorem]{Lemma}
\newtheorem{proposition}[theorem]{Proposition}
\newtheorem{corollary}[theorem]{Corollary}
\theoremstyle{definition}
\newtheorem{definition}[theorem]{Definition}
\theoremstyle{remark}
\newtheorem{remark}[theorem]{Remark}
\numberwithin{equation}{section}
\setlist{topsep=4pt,itemsep=2pt,parsep=0pt}
\setlength{\parskip}{2pt}
\emergencystretch=2em
\title{EXPTIME-completeness of group Weisfeiler--Leman\\
identification and equivalence}
\author{}
\date{}
\begin{document}
\maketitle
\vspace{-2em}
\begin{abstract}
Fix an odd prime $p$. A dense alternating tensor over $\mathbb F_p$ describes a
finite group whose order can be exponential in the length of its description.
Given the tensor and a dimension $k$, deciding whether $k$-dimensional
Weisfeiler--Leman refinement identifies this group among all finite groups
is EXPTIME-complete. Deciding equivalence of two tensor-described groups is
also EXPTIME-complete. Both statements hold for all three versions of group
refinement and for binary and unary $k$. The identification reduction outputs
groups of class two and exponent $p$ with center equal to their derived
subgroup. The equivalence reduction outputs nonisomorphic pairs on every
input. The proof encodes a bounded machine computation as an affine system.
Solvability determines whether a twist of a group bracket can be removed,
and local solutions give Duplicator's strategy below a prescribed dimension.
\end{abstract}

\section{The problems and the main theorem}
\label{sec:statement}
Fix an odd prime $p$ and write $\F=\mathbb F_p$.
For $V=\F^d$ and $W=\F^s$, an alternating bilinear map $b:V\times V\to W$
defines a group $G_b$ on $V\oplus W$ by
\begin{equation}
 (v,z)(v',z')=(v+v',z+z'+\tfrac12 b(v,v')).
 \label{eq:tensor-group}
\end{equation}
Bilinearity gives associativity. The identity is $(0,0)$, the inverse of
$(v,z)$ is $(-v,-z)$, and the commutator $ghg^{-1}h^{-1}$ of elements
with first coordinates $v,v'$ is $(0,b(v,v'))$. In particular
$|G_b|=p^{d+s}$. We use $\Alt(V,W)$ for the space of such maps.

A tensor description consists of $d,s$ followed by all $d^2s$
coefficients $b(e_i,e_j)_\ell$ in standard basis order. Each dimension
and each coefficient is written as $1^a0$, with coefficients represented
by integers $0\le a<p$. A valid array has zero diagonal and is
skew-symmetric. In the identification problem one tensor is followed by
$k\ge1$. In the equivalence problem two consecutive tensors are followed
by $k$. The suffix for $k$ is either its binary expansion beginning in
$1$, or the string $1^k$, according to the problem under consideration.
Malformed strings are rejected. In particular the input length charges
both dimensions even when $b=0$. The bases are used only for describing
the input. The groups given to refinement have no distinguished elements.

\paragraph{Colors and equivalence.}
For a $k$-tuple $\bar a$ in a finite set $S$, write
$\bar a[i\leftarrow x]$ for replacement of coordinate $i$ by $x$.
Starting from a color $c_0$, use the joint replacement rule
\[
 c_{t+1}(\bar a)=\left(c_t(\bar a), \left\{\!\!\left\{
 \big(c_t(\bar a[i\leftarrow x])\big)_{i=1}^k : x\in S
 \right\}\!\!\right\}\right),
\]
where double braces denote a multiset. The same $x$ is used in all
coordinates of a replacement vector. Colors are compared with common
names across the two structures.

In Version I, $S=G$ and $c_0$ records all equalities $a_i=a_j$ and
products $a_i a_j=a_\ell$. In Version II, $S=G$ and $c_0$ records the
group words $w$ for which $w(\bar a)=1$. Equality of Version-II initial
colors means that the marked subgroups are isomorphic under $a_i\mapsto
a_i'$, since they are quotients of the same free group by the same kernel.
In these versions the displayed rule applies also at $k=1$.

In Version III, form the uncolored simple graph $\Gamma(G)$ with one
vertex per group element and a fresh path $u_0u_1u_2u_3$ for each
ordered pair $(g,h)$. Join $u_0$ to $g$, $u_1$ to $h$, and $u_3$ to
$gh$. For $k\ge2$ apply the displayed rule with $S=V(\Gamma(G))$
and initial color recording equalities and adjacencies. For $k=1$ start
with the constant vertex color and refine by the old color together with
the multiset of neighbor colors. In both cases collect the final
histogram only on tuples of group vertices. These are the three group
versions of Brachter and Schweitzer~\cite{BS2020}, with the preceding
conventions fixing dimension one.

Write $G\equiv_k^J H$ if the multisets of colors of tuples in $G^k$ and
$H^k$ agree at every round in Version $J$. This can equivalently be
tested when the common partition stabilizes, since each update retains
the old color and the sets of tuples are finite. We say that $k$-WL
\emph{identifies} $G$ if
\begin{equation}
 \text{for every finite group }H,\qquad
 G\equiv_k^J H\quad\Longrightarrow\quad G\cong H .
 \label{eq:identifies}
\end{equation}
In particular, comparison groups in this definition may have no tensor
presentation of the input form.

\begin{theorem}\label{thm:main}
For every fixed odd prime $p$, every $J\in\{\mathrm I,\mathrm{II},\mathrm{III}\}$
and either encoding of $k$, the following two languages are EXPTIME-complete
under deterministic polynomial-time many-one reductions.
\begin{enumerate}[label=\textup{(\roman*)}]
\item The identification language consists of the inputs $(b,k)$ such
that $k$-WL identifies $G_b$ among all finite groups. Hardness holds
with every reduction output describing a group of nilpotency class
exactly two and exponent $p$ with $Z(G_b)=G_b'$.
\item The equivalence language consists of the inputs $(b,b',k)$ with
$G_b\equiv_k^J G_{b'}$. Every language in EXPTIME reduces to it with
$G_b\not\cong G_{b'}$ on every output, including yes-instances.
\end{enumerate}
\end{theorem}
Here EXPTIME consists of languages decided in time $2^{c(n+1)^e}$ for some constants $c,e$, where $n$ is the input length. A reduction writes the full target
word in polynomial time. Both definitions can use deterministic one-tape
machines with fixed finite alphabets. The proposition
\texttt{GroupWL.Statement.ManuscriptClaim} is the formal statement of
Theorem~\ref{thm:main}.

\begin{corollary}[Unbounded WL dimension]\label{cor:unbounded}
For every fixed odd prime $p$ and every $k\ge1$, there is one pair of
nonisomorphic finite groups $G,H$ of class two and exponent $p$ with
$Z(G)=G'$ and $Z(H)=H'$ such that $G\equiv_k^J H$ simultaneously for
all three versions.
\end{corollary}

\paragraph{Proof outline.}
An addressed circuit describes an exponentially long computation with
polynomially many connection rules (Appendix~\ref{app:compiler}).
Section~\ref{hs:section} converts these rules into an affine system
$Az=\beta$. It has a fixed inconsistent right-hand side that is locally
solvable whenever the circuit rejects. Sections~\ref{hyb:algebra}
and~\ref{hyb:residuals} turn this obstruction into a nonisomorphic pair
and convert its local solutions into winning bijections on whole groups.
On an accepting computation we must rule out \emph{every} nonisomorphic
comparison group. Four pebbled group elements recover its vector spaces
and its evaluation map up to a single alternating twist
(Section~\ref{hyb:recognition}). Decoding the twist yields a solvable
right-hand side, and hence an isomorphism (Section~\ref{sec:reduction}).
Section~\ref{hu:section} gives the decision algorithms.

\section{The affine computation source}
\label{hs:section}

Let $C$ be the circuit constructed in Lemma~\ref{lem:machine-schemas}
of Appendix~\ref{app:compiler}. Its addressed layout follows~\cite{OpenAI2026}.
Its gates are indexed by a finite type
and an address in $\{0,1\}^r$, with $r\ge32$ polynomial in the input
length. A connection is specified by its endpoint types, its target
port, a product domain and a coordinatewise injective bit map on that
domain. Seeds and tests also have product domains. A scalar gate takes
the OR of its seeds and incoming signals. A two-port AND gate takes
the conjunction of the ORs on its two ports. The circuit is acyclic
and accepts if any tested gate is true. For the fixed source machine
there are $O(1)$ gate types and $O(r^2)$ schemas. The affine system below will be listed from these schemas.

\paragraph{Comparisons and blocks.}
A scalar gate has one coordinate. An AND gate has two coordinate input
forms and their sum as output.  Add a scalar broadcast type \(\mathsf{Bcast}\) and
a dummy with zero form.  Compare each seed with the dummy's zero form.
Besides circuit connections, add connections from tests to \(\mathsf{Bcast}\), from
\(\mathsf{Bcast}\) to itself flipping each address bit, from \(\mathsf{Bcast}\) to every gate
coordinate, and identity self-connections on every gate coordinate and
on \(\mathsf{Bcast}\), with two copies for \(\mathsf{Bcast}\).
Replace every directed connection \(o\to s\) by a private two-coordinate
buffer and the three comparisons
\begin{equation}
 o-z_1=b_1,\qquad o-z_2=b_2,\qquad z_1+z_2-s=b_3.
 \label{hs:buffer}
\end{equation}
The first two keep the source address. The third uses the connection's
address map.  All three use its product domain, outside which the
buffer is unused.
Write \(d_P\in\{1,2\}\) for each real type's dimension and
\(\lambda_{e,0},\lambda_{e,1}\) for a comparison's endpoint forms.

Put \(K=r+1\).  Subdivide each comparison \(e:P\to Q\) into a wire
of \(4K+1\) links with private scalar internal types.
Write \(A_{t,i}=D_{e,i}\) when \(t\) is internal to wire \(e\), and
\(A_{t,i}=\{0,1\}\) when \(t\) is real.
The label matches \(M_{\ell,i}\) on a link of wire \(e\) are
\(\{(a,a):a\in D_{e,i}\}\), except on the last link, where they are
\(\{(a,\sigma_{e,i}(a)):a\in D_{e,i}\}\).
Link endpoint forms \(L_{\ell,t},L_{\ell,u}\) are
identity forms internally and the prescribed comparison forms at real
endpoints.  Parallel schemas and self-schemas are allowed.

The corners are \((t,i,a)\), with \(a\in A_{t,i}\).
For \(a\in A_{t,i}\), \(b\in A_{t,i+1}\), the horizontal readout
\(x_{i,t,a,b}\in\F^{d_t}\) is the edge
\((t,i,a)\to(t,i+1,b)\).
For \(\ell:t\to u\) and \((a,c)\in M_{\ell,i}\), the vertical readout
\(y_{i,\ell,a,c}\in\F\) is the edge
\((t,i,a)\to(u,i,c)\).
There is a site block \(0=\beta_B\) owning each horizontal readout.
For consecutive matches \((a,c)\in M_{\ell,i}\) and
\((b,d)\in M_{\ell,i+1}\), make the square equation
\begin{equation}
 L_{\ell,t}x_{i,t,a,b}-L_{\ell,u}x_{i,u,c,d}
       +y_{i+1,\ell,b,d}-y_{i,\ell,a,c}=\beta_B.
 \label{hs:square}
\end{equation}
Rows are indexed by $i=1,\ldots,r$ cyclically. We suppress label subscripts when the selected square determines them.
A square uses at most four readouts and five
scalar coordinates, since at least one endpoint is internal.
There are \(O(r^4)\) blocks, readouts, and incidences.
Together these equations form \(Az=\beta\). Site scopes are retained
even though their coefficient is zero. Every horizontal readout occurs at its site.
Every vertical readout belongs to a square because its label on the link extends
to an allowed label on the next row since each coordinate domain is
nonempty.

For a set \(S\) of blocks let \(J(S)\) be its readouts.
An \emph{exact local family through \(K\) blocks} consists of nonempty
sets \(\mathcal H_S\) of assignments on \(J(S)\) satisfying the equations
in \(S\), for \(|S|\leq K\), with
\(\mathcal H_T|_{J(S)}=\mathcal H_S\) whenever \(S\subseteq T\).
Let \(B(i,\ell,a)\) be the square selected by a full wire address
\(a\in D_e\), transmitted by \(\sigma_e\) at its last endpoint.  Define
\[
 \Delta_\ell(a)=\sum_i\beta_{B(i,\ell,a)},\qquad
 \Delta_e(a)=\sum_{\ell\text{ on }e}\Delta_\ell(a).
\]
Each is a \emph{row sum}, meaning a sum of functions of
\((a_i,a_{i+1})\).  The scalar-image conditions are
\begin{equation}
 \lambda_{e,0}(Q_{P,a})
 =\Delta_e(a)+\lambda_{e,1}(Q_{Q,\sigma_e(a)}),
 \qquad \varnothing\neq Q_{P,a}\subseteq\F^{d_P}.
 \label{hs:images}
\end{equation}
For each comparison this is equality of two sets of scalars. Different comparisons may use different vectors in the same set $Q_{P,a}$.

\begin{lemma}[Affine source]
\label{hs:source}
The construction has the following properties.
\begin{enumerate}
\item On an accepting input, every \(\beta\) that vanishes at sites
and admits the sets \eqref{hs:images} belongs to \(\operatorname{im}A\).
\item One explicitly computable shift \(\beta^*\) is outside
\(\operatorname{im}A\) on every input.  On a rejecting input it has
an exact local family through \(K=r+1\) blocks.
\end{enumerate}
\end{lemma}
\begin{proof}
We prove the three ingredients, including the signed extension
argument needed in odd characteristic.

\paragraph{Global lifting.}
A global solution yields row-sum functions
\(h_t(a)=\sum_i x_{i,t,a_i,a_{i+1}}\), and summing squares around rows
and along wires gives
\begin{equation}
 \lambda_{e,0}h_P(a)-\lambda_{e,1}h_Q(\sigma_e(a))=\Delta_e(a).
 \label{hs:global}
\end{equation}
Conversely, suppose row-sum functions at real types satisfy
\eqref{hs:global}.  At an internal type \(s\) on wire \(e\), set
\(h_s=\lambda_{e,0}h_P-\sum_{\ell\leq s}\Delta_\ell\).
Choose row decompositions once per type and assign their terms to the
horizontal readouts.  On each link put $g_i=\beta_{B(i,\ell,a)}-L_{\ell,t}x_i+
L_{\ell,u}x'_i$, using the two assigned horizontal readouts of the
selected square. By the link equation for $h_t,h_u$, these residuals
sum to zero at every product address.
They have the form \(g_i(u,v)=Y_{i+1}(v)-Y_i(u)\):
take rectangular differences in neighboring coordinates. Since $r\ge3$, only the term $g_i$ involves both of them, so its mixed difference is zero. Subtracting its values at a fixed base pair gives
\(g_i=c_i+\alpha_i(u)+\gamma_i(v)\), with both one-variable terms normalized to vanish at the base labels. Varying coordinate \(i\) then gives
\(\gamma_{i-1}=-\alpha_i\), and \(\sum_i c_i=0\).
Choose \(C_{i+1}-C_i=c_i\) cyclically and put \(Y_i=-\alpha_i+C_i\).
These \(Y_i\) give the vertical readouts.
Matches are indexed bijectively by source-domain labels, so the
assignment satisfies every square.  Thus \eqref{hs:global} with
row-sum functions is also sufficient for a global solution.

\paragraph{Acceptance.}
In \eqref{hs:buffer}, a singleton source image forces both buffer
coordinates, hence the buffer and destination image, to be singletons.
At a seeded address the comparison with the dummy's zero form has a
singleton image on the dummy side, so it fixes the seed's scalar image.  Induction through the acyclic circuit
makes every true output singleton. A true AND has both coordinates
singleton.  An accepting test therefore makes one broadcast address
singleton.  Bit flips connect the address cube, and broadcasting then
makes the gate and broadcast sets singletons, say \(\{h_P(a)\}\).
For an identity self-connection on a coordinate \(h\), write \(\Delta_1,\Delta_2,\Delta_3\) for the row sums on its three comparison wires. Their equations give
\[
 z_1=h-\Delta_1,\quad z_2=h-\Delta_2,\quad
 h=\Delta_1+\Delta_2+\Delta_3.
\]
The last identity uses \(2h-h=h\), not characteristic two.
Consequently all gate and broadcast coordinates are row sums.
Each buffer coordinate is \(o(h_P)-\Delta_j\) on its product domain.
Extend its row terms arbitrarily outside that domain.
Give the dummy value zero.  These row-sum functions satisfy
\eqref{hs:global}, so global lifting proves the first assertion.

\paragraph{The rejecting image sets.}
Set all comparison constants \(b_e\) to zero except the third
comparison of one broadcast self-connection, which receives one.
On each wire put \(b_\ell=b_e\) on its last link and zero elsewhere,
and set
\begin{equation}
 \beta^*_B=\delta_i b_\ell,\qquad
 \delta_r=1,\quad \delta_i=0\ (i<r),
 \label{hs:seam}
\end{equation}
with zero site shifts.  The two broadcast self-connections force
\(h_{\mathsf{Bcast}}=0\) and \(h_{\mathsf{Bcast}}=1\) in any global solution, on every input.

On rejection give a scalar gate \(\{0\}\) if true and \(\F\) otherwise.
At an AND gate use the product of \(\{0\}\) or \(\F\) according to
its port truths. Its output image is \(\{0\}\) exactly when both are
true.  Give \(\mathsf{Bcast}\) and the dummy the full set.
A singleton source of an unshifted connection has singleton-zero
destination, so use buffer \(\{(0,0)\}\).
For a full source and nonempty destination image \(T\), use
\(\{(u,v):u+v\in b_3+T\}\).
Both coordinate images are full and its sum image is \(b_3+T\).
Outside a buffer's active domain choose any nonempty set.
Thus \eqref{hs:images} holds with \(\Delta_e=b_e\).

\paragraph{Exact local projection.}
For a set $S$ of at most $K=r+1$ blocks, form the graph on corners using
the horizontal and vertical readouts of $S$ as edges.
A site contributes one edge and a square contributes its four boundary
edges. Each block thus has connected support. A link is active if some
square in $S$ uses it. On each wire choose one inactive link as a cut.
This is possible since a wire has $4K+1$ links and $S$ activates at
most $K$. Removing those links attaches every internal type to one of
the two real endpoints of its wire. Call that endpoint its root. For
a corner on the target arm of wire $e$ align its bit label at coordinate
$i$ by $\sigma_{e,i}$, and on the source arm use the original label.
At a real type use the original label. Along an active vertical edge
the root and aligned bit agree at both ends.
For a horizontal readout the root is constant and its two aligned
labels lie in consecutive rows. The two horizontal sides of a square
have the same pair of aligned labels.

Put marker $\delta_i$ on a forward horizontal edge in row $i$, its
negative on a backward traversal, and zero on a vertical edge. A
component wraps if a closed walk has nonzero marker sum in $\F$.
For such a walk let $f_i$ be its signed number of horizontal traversals
from row $i$ to row $i+1$ in $\F$. Flow conservation on row $i$ gives
$f_i=f_{i-1}$, so every $f_i$ equals the nonzero marker sum.
Thus a wrapping subcollection uses at least one block in every row.
If a wrapping set $R$ has $r$ blocks, there is exactly one in each row,
and the walk meets a horizontal edge of each. Connectedness of block
supports then puts every used corner in the wrapping component. In each
row only one ordered pair of aligned labels is available for horizontal
traversals. Applying flow conservation at each aligned label separately
shows that the ending label in row $i-1$ equals the starting label in
row $i$, since distinct labels would give nonzero net flow at one of them. Active vertical
edges preserve that label. Hence all used corners have a common root
$P$ and their aligned labels equal $a_i$ for a single address $a$.
If two wrapping $r$-subsets of $S$ exist they differ in at most one
block each. Their common blocks cover all but at most one row. These
cyclic adjacent pairs determine every bit, so both subsets specify the
same $(P,a)$ for the chosen cuts.

Choose one vector $h_P$ at each real type. If a wrapping $r$-subset of
$S$ specifies $(P,a)$, impose $h_P\in Q_{P,a}$. Such a choice exists
by nonemptiness and the preceding uniqueness. For each internal type
of wire $e:P\to Q$ use shift $h_t=\lambda_{e,0}(h_P)$ on the source
arm and $h_t=b_e+\lambda_{e,1}(h_Q)$ on the target arm. Thus on every
active link $\ell:t\to u$,
$L_{\ell,t}h_t-L_{\ell,u}h_u=b_\ell$.
At each corner $v$ choose a potential $M_v$ in its type space.
On the oriented edges put
\begin{equation}
 x_{vw}=M_v-M_w+\delta_i h_t,\qquad
 y_{vw}=L_{\ell,t}M_v-L_{\ell,u}M_w .
 \label{hs:potentials}
\end{equation}
Let $\mathcal H_S$ consist of the assignments on $J(S)$ so obtained,
allowing every admissible choice of cuts, shifts and potentials.
Substitution into \eqref{hs:square} cancels the four potential terms
and gives $\delta_i b_\ell$ as required. Site equations also hold.
For $S\subseteq T$ every admissible choice for $T$ is admissible for
$S$, proving restriction.

For extension assume $S\subsetneq T$ and $|T|\le r+1$. Retain witnesses
for the given assignment on $J(S)$ and choose cuts for $T$. If $S$
does not wrap, take any admissible shifts for $T$. Otherwise
$|S|=r$, and the used corners of $S$ belong to one wrapping component.
If it meets a real type $P$, its root for either set of cuts is $P$.
On each active link the transport of a bit label from either endpoint
to the other is fixed, independent of the cut. Thus transport from any
old corner to $P$ along an old path gives the same coordinate
permutation for both cut choices. Consequently the old address is still
the address prescribed by $S$. Keep $h_P$. The shifts on all its used
wire arms are determined by this vector and the active link equations.
Here and below, $S$ already specifies the sole root/address requirement
of $T$, by the common-block argument above.

If the old component contains only internal types, it lies on a single
wire $e:P\to Q$. All its used internal nodes have the same shift $z$
and use the same original label coordinates. Since every row is
represented, they determine an address $a\in D_e$. Each bit occurs
as an allowed internal label. According to the old cut, either
$z\in\lambda_{e,0}(Q_{P,a})$ or
$z\in b_e+\lambda_{e,1}(Q_{Q,\sigma_e(a)})$.
These two sets coincide by \eqref{hs:images} with $\Delta_e=b_e$.
For the new cut, choose a vector in the appropriate $Q$-set
whose indicated affine image is $z$. This preserves $z$ at every old
used type, including when both endpoints have the same real type.
Choose the other real shifts freely. This constructs admissible shifts
on $T$ preserving the shifts of every old wrapping component.

On an old nonwrapping component the marker on $v\to w$ is $s_w-s_v$,
where $s_v$ is the integral from a fixed corner to $v$. Zero integrals
on closed walks give path independence. Put
$g_t=h_t^{\rm new}-h_t^{\rm old}$ and replace
$M_v$ by $M_v+s_vg_t$ on that component. The change on a horizontal
edge is $(s_v-s_w+\delta_i)g_t=0$. On an old vertical edge $s_v=s_w$
and $L_{\ell,t}g_t=L_{\ell,u}g_u$, so its value also stays fixed.
Keep the old potentials on wrapping components. All these changes refer
to disjoint sets of old corners, even if new edges merge components.
Choose arbitrary potentials at previously unused corners. This extends
the prescribed assignment to $J(T)$ and proves exact projection.
\end{proof}


\section{An affine system as a commutator obstruction}
\label{hyb:algebra}

The affine source has readouts $z_j\in V_j$, equations indexed by $I$,
and scopes $T_i\subseteq J$:
\begin{equation}
 (Az)_i=\sum_{j\in T_i}\lambda_{ij}(z_j)=\beta_i.
 \label{hyb:source-equations}
\end{equation}
Scope membership is retained even when $\lambda_{ij}=0$, and every
readout occurs in a scope.  Assume $I\ne\varnothing$, as is the case for every constructed instance.  All spaces and maps below are over $\F$.

\subsection{The tensor and its padding}

Use the bipartite incidence graph on $I\sqcup J$.  At its vertices put
\[
 N_i=\F\quad(i\in I),\qquad
 N_j=\left\{(h_i)_{i:j\in T_i}:
                 \sum_{i:j\in T_i}h_i\lambda_{ij}=0\text{ in }V_j^*\right\},
 \qquad R_v=\F\oplus N_v,
\]
with multiplication $(a,u)(b,w)=(ab,aw+bu)$.  Thus $N_v^2=0$, and an
element of $R_v$ is a unit exactly when its scalar part is nonzero.
For an incidence $ij$ let $R_{ij}=\F\oplus\F$ with the same
multiplication, and use the algebra maps
\[
 \phi_{ij}:R_i\to R_{ij},\quad(a,h)\mapsto(a,h),\qquad
 \psi_{ij}:R_j\to R_{ij},\quad(a,(h_t)_t)\mapsto(a,h_i).
\]
Choose $q\ge4$ and set
\[
 X^\circ=\bigoplus_v R_v^q,\qquad
 Z^\circ=\bigoplus_{v,\,a<b}R_v\ \oplus\bigoplus_{ij,\,a,b}R_{ij}.
\]
The alternating tensor $\omega:X^\circ\times X^\circ\to Z^\circ$ is
\begin{align}
 \omega(x,x')_{v,ab}&=x_{va}x'_{vb}-x'_{va}x_{vb},\notag\\
 \omega(x,x')_{ij,ab}
   &=\phi_{ij}(x_{ia})\psi_{ij}(x'_{jb})
       -\phi_{ij}(x'_{ia})\psi_{ij}(x_{jb}).
 \label{hyb:omega}
\end{align}
Let $L\le X^\circ$ consist of vectors whose scalar part in slot $1$
vanishes at every vertex. Its basis is obtained by deleting the
corresponding scalar units. In particular $L$ contains every vector
whose coordinates are radical elements.

We include the padding needed for the recognition argument from the start:
take a copy $D\cong Z^\circ$ and a copy $W_0\cong X^\circ\oplus D$, and put
\[
 X=X^\circ\oplus D,\qquad Z=Z^\circ\oplus W_0.
\]
Extend $\omega$ by zero on arguments in $D$ and with zero $W_0$ output.
In each $R_v$ use $1$ followed by a basis of $N_v$, using the same
basis in every copy of $R_v$. Use the induced coordinate bases on
the direct sums, using the field basis $1$ for the one-dimensional
equation and edge radicals.
For a readout radical use the kernel basis computed in
Section~\ref{sec:gaussian-bases}. All coordinate dot products, including
the Frobenius pairing $\langle B,C\rangle=\sum_{z,x} B_{zx}C_{zx}$,
refer to these fixed bases.  Define
\begin{equation}
 C=\{M_\omega(a):a\in L\},\qquad M_\omega(a)x=\omega(a,x),
 \qquad Y=C^\perp\le\Hom(X,Z).
 \label{hyb:C}
\end{equation}
Write $C^\circ$ for the restrictions of $C$ to $X^\circ$ with output
$Z^\circ$, and put $Y^\circ=(C^\circ)^\perp$.  Then
$Y=Y^\circ\oplus\Hom(D,Z^\circ)\oplus\Hom(X,W_0)$.
The map $b(x,d)=d$ onto $Z^\circ$ is an explicit
element of $Y$.  The free output summand contains the identity and a
cyclic shift on $X$, as required for the recognition frame.

For a linear map $\alpha:X\to Y$ write
\[
 \partial\alpha(x,x')=\alpha(x)x'-\alpha(x')x.
\]
On $X\oplus Y\oplus Z$ define the group $G_\varepsilon$ by
\begin{equation}
 (x,y,z)(x',y',z')=
 \left(x+x',y+y',z+z'+\tfrac12
       \bigl(y'x-yx'+\varepsilon(x,x')\bigr)\right).
 \label{hyb:group}
\end{equation}
It has class at most two and exponent $p$.  A shear
$(x,y,z)\mapsto(x,y+\alpha(x),z)$ maps $G_0$ to $G_\varepsilon$
multiplicatively exactly when $\partial\alpha=\varepsilon$.

\begin{lemma}[Rank and local multipliers]
\label{hyb:local}
For every nonzero $a\in X^\circ$, the map $M_\omega(a)$ has rank at least $q-1$.
Consequently $Yx=Z$ for every nonzero $x\in X$, and
$Z(G_\varepsilon)=G_\varepsilon'=Z$.

More generally, let $U\le X$ and let an alternating $\gamma:U\times U\to Z$
satisfy
\begin{equation}
 \gamma(u,v)=\omega(Pu,v)\quad(u,v\in U)
 \label{hyb:partial-symbol}
\end{equation}
for some linear $P:U\to L$.  Call a vertex $v$ good if its $q$ scalar
coordinate functionals have rank at least two on $U$.  At a good vertex
there is a unique $r_v\in R_v$ such that
\[
 (Pu)_v=r_vu_v\quad(u\in U).
\]
If $i,j$ are both good and incident, then
$\phi_{ij}(r_i)=\psi_{ij}(r_j)$.
\end{lemma}

\begin{proof}
For the rank assertion choose $a_{va}\ne0$ and evaluate on the $q-1$
unit vectors in the other coordinates of $R_v^q$.  In the internal
output $(v;a,b)$ their linear combination has coefficient
$\pm a_{va}t_b$, so these columns are independent.  Thus $C$ has no
nonzero rank-one matrix.  Failure of $Yx=Z$ would give precisely such a
matrix in $Y^\perp=C$, by pairing evaluation at $x$ with a nonzero
functional on $Z$.  The evaluation brackets span $Z$.
A central element must then have $x=0$, and its
$Y$-coordinate must vanish as an operator. Hence the center is $Z$.

For the local assertion, $\gamma(u,u)=0$ gives
$(Pu)_v\wedge u_v=0$ in all internal coordinates.  If the scalar
projection of $u_v$ is nonzero, one coordinate is a unit, so
$(Pu)_v=r_u u_v$ for a unique $r_u\in R_v$.  Choose two vectors whose
scalar projections are independent.  Their wedge has a unit coordinate.
Polarizing the displayed identity forces their multipliers to agree.
Every nonzero scalar projection is independent of at least one of these
two, so its multiplier is the same $r_v$.  For zero scalar projection,
add one of the chosen vectors and subtract again, using linearity of $P$.
This also covers $u_v=0$.
For an incidence with good endpoints choose $u\in U$ with a selected
scalar coordinate nonzero at each endpoint. Two proper kernels cannot
cover $U$.  Its diagonal incidence identity is
\[
 (\phi_{ij}(r_i)-\psi_{ij}(r_j))
       \phi_{ij}(u_{ia})\psi_{ij}(u_{jb})=0.
\]
The last two factors are units, proving the claim.
\end{proof}

Let $\tau_i$ be the radical basis element of $R_i=\F\oplus\F$.
Define $\varepsilon_\beta$ to have only the output $(i;1,2)$, on the
two scalar unit inputs at vertex $i$, with value $\beta_i\tau_i$.
Extend by alternation and bilinearity and set all other basis values to zero.

\begin{proposition}[The full obstruction quotient]
\label{hyb:cokernel}
Every alternating twist is congruent modulo $\im\partial$ to some
$\varepsilon_\beta$, and
\[
 \varepsilon_\beta\in\im\partial
       \quad\Longleftrightarrow\quad \beta\in\im A.
\]
These assertions include the free input and output padding above.
\end{proposition}

\begin{proof}
Pair alternating tensors by their values on unordered coordinate-basis
pairs.  Antisymmetrization gives
\begin{equation}
 \gamma\in(\im\partial)^\perp
 \quad\Longleftrightarrow\quad M_\gamma(x)\in C\quad(x\in X).
 \label{hyb:global-dual}
\end{equation}
Indeed, on a coordinate basis $(e_a)$ of $X$ the pairing with
$\partial\alpha$ equals
$\sum_a\langle M_\gamma(e_a),\alpha(e_a)\rangle$.
The values $\alpha(e_a)\in Y$ vary independently. The rank lemma
makes $M_\omega$ injective on $L$, so such a tensor has
the form \eqref{hyb:partial-symbol} on all of $X$.  All vertices are good.
Lemma~\ref{hyb:local} gives local multipliers at every vertex.
Since $P(X)\subseteq L$, testing on the scalar unit in slot $1$ at
each vertex forces the scalar part of its multiplier to vanish.  The remaining multipliers are
\[
 r_i=(0,h_i),\qquad r_j=(0,(h_i)_{i:j\in T_i}),\qquad h\in\ker A^*.
\]
Conversely every such $h$ defines a tensor $\gamma_h=\omega(r\,\cdot,\cdot)$
satisfying \eqref{hyb:global-dual}.  There are no extra padded terms:
every slice in $C$ is zero on $D$ and in output $W_0$, and alternation
also kills $D$ in the first argument.  Thus
\[
 (\im\partial)^\perp=\{\gamma_h:h\in\ker A^*\},\qquad
 \langle\gamma_h,\varepsilon_\beta\rangle=\sum_i h_i\beta_i.
\]
The last pairing vanishes for all $h\in\ker A^*$ exactly when
$\beta\in(\ker A^*)^\perp=\im A$. For an arbitrary twist $\varepsilon$,
the map $h\mapsto\langle\gamma_h,\varepsilon\rangle$ is a linear
functional on $\ker A^*$. Extend it to $\F^I$ and write the extension
as $h\mapsto\sum_i h_i\beta_i$. Then
$\varepsilon-\varepsilon_\beta$ annihilates $(\im\partial)^\perp$
and therefore lies in $\im\partial$.
\end{proof}

\subsection{Computing the tensor coefficients}
\label{sec:gaussian-bases}

We give one procedure for both kernels used by the reduction.
For an $m\times n$ matrix over $\F$, start with the $n\times n$
identity matrix $B$. Process the rows in reverse order. The columns
of $B$ form a basis of the kernel of the rows processed so far.
For the next row $a$, compute $c=aB$. If $c=0$, retain $B$.
Otherwise, swap its first nonzero coordinate with its first coordinate,
permuting the columns of $B$ in the same way. If the current number of
columns is $d$, take as the new columns
\[
 B_j-(c_j/c_1)B_1,\qquad 2\le j\le d,
\]
where $B_j$ and $c_j$ denote the columns and entries after the swap.
These columns span exactly the vectors $Bt$ with $ct=0$, since the
equation determines $t_1$ from the other coordinates. The uniqueness
of those other coordinates proves independence. This also handles
zero rows, zero-dimensional kernels and matrices with no rows.
Storing the current matrix explicitly costs $O((n+1)^2)$ field
entries, and all row updates together use
$O((m+1)(n+1)^2)$ field operations.

At a readout $j$, order the incident equations and let $e_b$ range over
the given coordinate basis of $V_j$. Apply this procedure to the matrix
with entries $\lambda_{ij}(e_b)$ in row $b$ and column $i$. Its kernel
is exactly $N_j$. Store its basis columns in the original incidence
coordinates. In this basis $\psi_{ij}$ reads the $i$-th entry of a
stored column on a radical basis vector. Internal multiplication uses
$1\cdot u=u$ and $N_v^2=0$. Thus \eqref{hyb:omega} computes the
coordinates of $\omega$ directly, on every pair of coordinate units.
The scalar restrictions defining $L$ select coordinate positions.
Applying the same procedure to these restrictions, or deleting
the restricted units, gives a basis $(a_h)$ of $L$.

Form a matrix with one row for each $a_h$ and with columns indexed by
$(z,x)$ in $Z\times X$ coordinate order. Its entries are
\[
           H_{h,(z,x)}=\omega(a_h,e_x)_z .
\]
Here and below a coordinate-order index ranges over basis positions.
Equation~\eqref{hyb:C} says that a list of matrix entries belongs
to $Y$ precisely when it is killed by $H$. Apply the procedure to
$H$ and view its basis columns as operators $y_t:X\to Z$.
The basis of the group input is now the ordered concatenation
of the coordinate units $e_x$ of $X$ and these $y_t$.
In this basis the entire alternating tensor in \eqref{hyb:group} is
\begin{equation}
 \begin{aligned}
 b_\varepsilon(e_x,e_{x'})&=\varepsilon(e_x,e_{x'}),&
 b_\varepsilon(y_s,y_t)&=0,\\
 b_\varepsilon(e_x,y_t)_z&=(y_t)_{zx},&
 b_\varepsilon(y_t,e_x)_z&=-(y_t)_{zx}.
 \end{aligned}
 \label{eq:computed-bracket}
\end{equation}
For the two reduction outputs insert $\varepsilon=0$ and
$\varepsilon=\varepsilon_{\beta^*}$, respectively. Both use the same
matrices $H$ and $y_t$.

\section{Coherent residuals for arbitrary group queries}
\label{hyb:residuals}

Suppose the affine instance has nonempty exact families
$\mathcal H_S$ for all $S\subseteq I$ with $|S|\le K$:
they assign the readouts in $J(S)=\bigcup_{i\in S}T_i$, satisfy the
equations in $S$, and
$\mathcal H_T|_{J(S)}=\mathcal H_S$ whenever $S\subseteq T$.

For any retained subspace $E\le X$ set
\[
 \begin{aligned}
 U&=E^\perp,& Y_E&=\{y\in Y:y(E)=0\},\\
 \mathcal R_E&=\{\rho\in\Alt(X,Z):\rho(E,X)=0\},&
 \mathcal B_E&=\{b\in\Hom(X,Y_E):b(E)=0\}.
 \end{aligned}
\]
Then $\partial$ maps $\mathcal B_E$ into $\mathcal R_E$.
Its exact relative dual criterion is
\begin{equation}
 \gamma\in(\partial\mathcal B_E)^\perp
 \quad\Longleftrightarrow\quad
 \gamma(u,v)=\omega(Pu,v)\ (u,v\in U)
 \text{ for some }P:U\to L.
 \label{hyb:relative-dual}
\end{equation}
To see this, identify \((X/E)^*\) with \(U\) using the coordinate
pairing on \(X\). Evaluation gives perfect pairings
between \(\Hom(X/E,Z)\) and \(\Hom(U,Z)\), and between
\(\Alt(X/E,Z)\) and \(\Alt(U,Z)\), using the dot product on \(Z\).
The second pairing uses exterior powers, so it makes no choice of a
complement to \(E\). The inclusion
\(\Hom(X/E,Z)\hookrightarrow\Hom(X,Z)\) is adjoint to restriction to
\(U\). Since \(Y_E\) is the inverse image of \(Y=C^\perp\) under this
inclusion, its annihilator is \(C|_U\). For \(u\in U\) and \(y\in Y_E\)
the map \(b(x)=\langle u,x\rangle y\) belongs to \(\mathcal B_E\).
These rank-one maps span it. Expanding the alternating difference gives
\[
 \langle\gamma,\partial b\rangle
   =\langle\gamma(u,\cdot),y\rangle .
\]
Indeed the two terms on each unordered pair collect all off-diagonal
terms in the contraction on the right, while the diagonal is zero.
Thus the contractions \(\gamma(u,\cdot)\) must all lie in \(C|_U\).
The map \(L\to C|_U\), \(a\mapsto\omega(a,\cdot)|_U\), is surjective.
Compose a linear section with the contraction map to obtain \(P\).
The same calculation proves the converse in~\eqref{hyb:relative-dual}.

Call a vertex good at $E$ according to Lemma~\ref{hyb:local} with
$U=E^\perp$, and bad otherwise.  If $S_v$ is the span of its $q$ scalar
unit vectors, the rank of those scalar functionals on $U$ is
$q-\dim(E\cap S_v)$.  The $S_v$ have disjoint support, so
\begin{equation}
 (q-1)\#\{v\text{ bad at }E\}\le\dim E.
 \label{hyb:bad-budget}
\end{equation}
Choose an owner $o(j)\in I$ with $j\in T_{o(j)}$ for each readout.
Let $O(E)$ contain every bad equation and the owner of every bad
readout.  It is monotone in $E$, and $|O(E)|\le\lfloor\dim E/(q-1)\rfloor$.
A total readout assignment $z$ is admissible at $E$ if
$z|_{J(O(E))}\in\mathcal H_{O(E)}$.  Its residual
$s_i(z)=\beta_i-(Az)_i$ is zero at every bad equation.

For each good equation $i$, choose $u,v\in U$ and coordinates $a<b$
whose local minor $t=u_{ia}v_{ib}-v_{ia}u_{ib}\in R_i$ is a unit.
Choose $z_0\in R_i$ representing, under the coordinate pairing, the
functional $w\mapsto\operatorname{rad}(t^{-1}w)$.  In the internal
output $(i;a,b)$ put
\[
 \eta_{E,i}(x,x')=
   \bigl(\langle u,x\rangle\langle v,x'\rangle
          -\langle v,x\rangle\langle u,x'\rangle\bigr)z_0,
\]
and set its other outputs to zero.  It belongs to $\mathcal R_E$.
For a relative obstruction in \eqref{hyb:relative-dual}, with local
multiplier $r_i=(a_i,h_i)$, the natural dual pairing gives
\begin{equation}
 \langle\gamma,\eta_{E,i}\rangle
     =\operatorname{rad}(t^{-1}r_it)=h_i.
 \label{hyb:normalized-residual}
\end{equation}
In particular this value is independent of the chosen minor.
Fix these choices once for each $E$, and define
\begin{equation}
 \rho_E(z)=\sum_{i\text{ good at }E}s_i(z)\eta_{E,i}.
 \label{hyb:rho}
\end{equation}
Its pairing with every global $\gamma_h$ equals
$\sum_i h_i(\beta_i-(Az)_i)=\sum_i h_i\beta_i$.
Bad equations contribute zero.  Proposition~\ref{hyb:cokernel} therefore gives
\begin{equation}
 \rho_E(z)\in\mathcal R_E,\qquad
 \varepsilon_\beta-\rho_E(z)\in\im\partial.
 \label{hyb:global-residual}
\end{equation}

\begin{lemma}[Coherent residuals]
\label{hyb:coherence}
Suppose $E\subseteq F$, $|O(F)|\le K$, and admissible $z_E,z_F$ agree
on $J(O(E))$.  Then
$\rho_F(z_F)-\rho_E(z_E)\in\partial\mathcal B_E$.
\end{lemma}

\begin{proof}
Both tensors belong to $\mathcal R_E$.  Pair their difference with a
relative obstruction \eqref{hyb:relative-dual}.  At an equation that
stays good, both normalizations read the same multiplier, by restricting
$P$ to $F^\perp\subseteq E^\perp$.  At one that becomes bad, its new
residual is zero.  The pairing is consequently
\[
 -\sum_{i\text{ good at }E}\sum_{j\in T_i}
       h_i\lambda_{ij}(z_{F,j}-z_{E,j}).
\]
If a readout changes, it is good at $E$, because otherwise its owner
pins it.  Every adjacent equation is good as well, because a bad
equation pins its whole scope.  On this complete good star,
Lemma~\ref{hyb:local} identifies the coefficients $h_i$ with the radical
coordinates of $r_j\in R_j$. Their definition gives
$\sum_{i:j\in T_i}h_i\lambda_{ij}=0$.
Grouping by changed readouts makes the pairing zero.
Equation~\eqref{hyb:relative-dual} proves the assertion.
\end{proof}

\begin{theorem}[Exact shear families and the bijective strategy]
\label{hyb:forward}
If the affine instance has the exact families above and
$m<(q-1)(K+1)$, Duplicator wins the bijective $m$-pebble game on
$G_0,G_{\varepsilon_\beta}$ with marked-subgroup isomorphism as its
atomic condition.  This includes arbitrary elements of all padded
input, matrix, and output sectors.
\end{theorem}

\begin{proof}
For each $E$ with $|O(E)|\le K$, let $\mathcal P_E$ consist of
restrictions $\alpha|_E$ of global maps satisfying
\[
 \partial\alpha=\varepsilon_\beta-\rho_E(z)
       \quad\text{for some assignment }z\text{ admissible at }E.
\]
These sets are nonempty by \eqref{hyb:global-residual}.  They have exact
restriction maps.  To extend a witness at $E\subseteq F$, first extend
its source assignment on $J(O(E))$ to one on $J(O(F))$.  Complete it
arbitrarily to a total $z_F$.  Lemma~\ref{hyb:coherence} supplies
$b\in\mathcal B_E$ with
$\partial b=\rho_F(z_F)-\rho_E(z_E)$.
Then $\alpha-b$ is a witness at $F$ and agrees with $\alpha$ on $E$.
Conversely, an admissible assignment at $F$ is admissible at $E$.
Apply the same lemma with this assignment held fixed and add the
correction to a witness at $F$.  Its restriction belongs to
$\mathcal P_E$ because the correction vanishes on $E$.

Maintain that the pebbled correspondence is the shear of some
$a\in\mathcal P_E$, where $E$ spans the retained $X$-coordinates.
Such a shear is an isomorphism on the full subgroup
$E\oplus Y\oplus Z$, since its residual vanishes on $E\times X$.
After removing a pebble, for every $x\in X$ independently choose an
extension $a_x\in\mathcal P_{E+\F x}$ of the retained $a$, and announce
the single bijection
\[
 (x,y,z)\longmapsto(x,y+a_x(x),z).
\]
It is a translation on each fixed $x,z$ fiber.  After the selected
move, its corresponding $a_x$ witnesses the invariant for all
retained pebbles.  Selection on the other side uses the inverse
bijection.  Finally \eqref{hyb:bad-budget} gives $|O(E)|\le K$ for
every span of at most $m$ query directions.  No restriction on the
support of those directions, including their $D$-coordinates, has
been used.
\end{proof}


\section{A four-anchor frame for every comparison group}
\label{hyb:recognition}

Recall the padded spaces
\(X=X^\circ\oplus D\), \(Z=Z^\circ\oplus W_0\), and
\(Y=Y^\circ\oplus\Hom(D,Z^\circ)\oplus\Hom(X,W_0)\), where
\(D\cong Z^\circ\) and \(W_0\cong X\).
The group is \eqref{hyb:group}, and Lemma~\ref{hyb:local} gives
\(Yx=Z\) for every nonzero \(x\in X\).
No named coordinates are added to the input group.

We use the Version-II bijective game of Lemma~\ref{lem:paper-wl-game}, with the marked-subgroup test
made after pickup.  A winning position with \(k+1\) pairs preserves
counting formulas in the group-word language using at most \(k\)
variable symbols, including parameters.  Here is the needed elementary
argument.  
Use induction on the formula with pairs retained on the free variables.
For an atomic equality, pick up a reserve pair: the marked-subgroup test
preserves every word equality in the retained elements. Negation and
conjunction preserve agreement by induction. For a counting quantifier
on a variable \(v\), release its old pair, if present. A winning
bijection on the two groups matches extensions that are winning
positions. Induction applied with each possible new value of \(v\)
shows that it matches the two sets of satisfying values bijectively.
Therefore the counting quantifier has the same truth value. The same
argument with no free variables transfers sentences. Formula length
and quantifier depth do not enter the pair count.
\begin{theorem}[Recognition and the zero-twist criterion]
\label{hyb:allmates}
If Duplicator wins the Version-II game with at least \(40\) pairs
between \(G_0\) and an arbitrary finite group \(H\), then
\[
 H\cong G_\eta
 \quad\text{for some }\eta\in\Alt(X,Z).
\]
The spaces and evaluation map in this conclusion are exactly the
prescribed \(X,Y,Z\).  Four retained anchors name the classes of the two successive quotients
and each central element.  Moreover, for every \(\eta\),
\begin{equation}
 G_0\cong G_\eta
 \quad\Longleftrightarrow\quad
 \eta\in\operatorname{im}\partial.
 \label{hyb:zero-isomorphism}
\end{equation}
\end{theorem}
\begin{proof}
Put \(n=\dim X\), \(m=\dim Z\),
\(s=\dim Y\), and \(T=p^{n+s}\).
Lemma~\ref{hyb:local} gives
\(Yx=Z\) for every \(x\neq0\).

For \(g=(x,y,z)\), its centralizer is the kernel of the linear map
\((x',y',z')\mapsto y'x-yx'+\eta(x,x')\). If \(x\ne0\),
the image is all of \(Z\), giving centralizer size precisely
\(p^{n+s}=T\). If \(x=0\), the rank is at most \(n\) and
\(m=n+\dim Z^\circ>n\). Its centralizer is strictly larger than \(T\).
Consequently
\[
 \mathcal B=Y\oplus Z=\{g:|C(g)|>T\},\qquad Z(G_\eta)=Z.
\]
Use the two-variable predicates
\[
 \mathsf B(u):=\exists^{\geq T+1}v\ [u,v]=1,
 \qquad \mathsf Z(u):=\forall v\ [u,v]=1.
\]
Exponent \(p\), class at most two, the assertion that \(\mathsf B\)
is an abelian subgroup, and the sizes of \(H,\mathsf B(H),\mathsf Z(H)\)
transfer.  Hence, writing \(\mathcal B_H=\mathsf B(H)\),
\begin{equation}
 |H/\mathcal B_H|=p^n,\qquad
 |\mathcal B_H/Z(H)|=p^s,\qquad |Z(H)|=p^m.
 \label{hyb:sizes}
\end{equation}
It contains \(Z(H)\) and is normal because its defining
centralizer-size condition is conjugation invariant.

Choose a basis \(e_0,\ldots,e_{n-1}\) of \(X\), with
\(e_0\in D\), and let \(Ae_i=e_{i+1\bmod n}\).
Pebble the four elements
\[
 a_0=I:X\to W_0,\qquad
 a_1=A:X\to W_0,\qquad
 b(x,d)=d\in Z^\circ,\qquad t=(e_0,0,0).
\]
The first three belong to \(Y\) and have all other
coordinates zero.  Fix the four winning responses as parameters in
\(H\). The first three lie in \(\mathcal B_H\) by transfer of
\(\mathsf B\).
The relations
\[
 [u,a_0]=1\ \Longleftrightarrow\ \mathsf B(u),
 \qquad \forall u\exists v\ [v,a_0]=[u,a_1]
\]
transfer.  Thus commutation with \(a_0\) embeds \(H/\mathcal B_H\)
in its center, and \(a_1\) defines an endomorphism \(A_H\) of that
quotient.

For \(x=\sum_{i<n}x_i e_i\), name its quotient class by a formula
\(\delta_x(u)\).  Starting with \(\delta_0(u)=\mathsf B(u)\),
Horner evaluation implements the polynomial step \(f\mapsto Uf+c\) by
\begin{equation}
 \delta_{Uf+c}(u):=
 \exists v\bigl(\delta_f(v)\land
       [ut^{-c},a_0]=[v,a_1]\bigr).
 \label{hyb:horner}
\end{equation}
Only the working names \(u,v\) are used: they alternate in the
recursive subformula, where the other outer name is not free.
These formulas name the cosets \(\sum_i x_iA_H^it\mathcal B_H\).
Their pairwise disjointness transfers from the reference group.
Equation~\eqref{hyb:sizes} then makes them all quotient classes,
identified linearly with \(X\).

Let \(\iota:Z^\circ\to D\leq X\) be the chosen copy map.
For \(z=(z_1,z_0)\in Z^\circ\oplus W_0\), using the fixed copy
\(W_0\cong X\), define
\begin{equation}
 \Phi_z(w):=\exists u\exists v\bigl(
   \delta_{z_0}(u)\land\delta_{\iota(z_1)}(v)
   \land w=[u,a_0][v,b]\bigr).
 \label{hyb:central-name}
\end{equation}
Since \(\mathcal B_H\) is abelian, these commutators depend only on
the named quotient classes.  Each formula therefore names one central
element, linearly in \(z\).  Transferred pairwise disjointness and
\eqref{hyb:sizes} identify all of \(Z(H)\) with \(Z\).

For each \(y\in Y\), set
\begin{equation}
 \Psi_y(w):=\mathsf B(w)\land
 \bigwedge_{i<n}\forall u\bigl(
     \delta_{e_i}(u)\Longrightarrow
     \Phi_{y(e_i)}([u,w])\bigr).
 \label{hyb:matrix-name}
\end{equation}
Every sentence \(\exists w\,\Psi_y(w)\) transfers.  Different
\(y\)'s give different central cosets.  Two realizations of the same
formula differ by an element commuting with \(\mathcal B_H\) and
with a basis of \(H/\mathcal B_H\), hence by a central element.
The middle count in \eqref{hyb:sizes} makes these all the classes of
\(\mathcal B_H/Z(H)\).  Commutation is linear, so this identifies
that quotient linearly with exactly \(Y\), with exactly
its original evaluation map.

Here is explicit variable accounting.  In the last two rows every
one of the four anchor names is reserved. The working count includes
all free query names.  Existential closure and conjunctions reuse the
same names.
\begin{center}
\begin{tabular}{lc}
\hline
Formula or assertion & Working names \\
\hline
\(\mathsf B,\mathsf Z\) & 2 \\
Subgroup, exponent, class and cardinality assertions & at most 3 \\
\(\delta_x\), including its base counting quantifier & 2 \\
\(\Phi_z(w)\) & 4 \\
\(\Psi_y(w)\), or \(\Phi_z([u,v])\) & 5 \\
\hline
\end{tabular}
\end{center}
For example, \(\Phi\) uses its output, two representatives and one
Horner helper. Substituting a two-query commutator for the output adds
one name.  Thus at most nine variable symbols occur, leaving the
required pickup reserve well within forty pairs.


To recover the group law, put \(d=2^{-1}\bmod p\),
use \([g,h]=ghg^{-1}h^{-1}\), and define
\(g\oplus h=gh[g,h]^{-d}\). Commutators are central and multiply
in either argument. Thus both bracketings of \(g\oplus h\oplus k\)
equal \(ghk([g,h][g,k][h,k])^{-d}\). Also
\(hg[h,g]^{-d}=gh[g,h]^{d-1}=gh[g,h]^{-d}\).
The identity and inverses agree with the group ones, and repeated
addition of \(g\) gives its powers. The resulting abelian group has
exponent \(p\) and hence is an \(\F\)-vector space. Its commutator
is alternating and \(\F\)-bilinear, since commutators with central
elements vanish. The original multiplication is recovered as
\(gh=g\oplus h\oplus d[g,h]\), with central elements written additively
in the last term.

In this vector space the subgroups \(Z(H)\subseteq\mathcal B_H\)
are subspaces. Choose linear sections of \(H\to H/\mathcal B_H=X\)
and \(\mathcal B_H\to\mathcal B_H/Z(H)=Y\).
The bracket on the second section is zero, and its bracket with the
first has the evaluation just recovered. The bracket on two lifts of
\(X\) defines an alternating \(\eta:X^2\to Z\).
The multiplication in the resulting coordinates is precisely
\eqref{hyb:group}, yielding \(H\cong G_\eta\).
In these coordinates the four matched anchors have the prescribed
quotient label of the seed and matrix labels of the other three:
the latter follow by transferring
\(\Psi_{a_0}(a_0),\Psi_{a_1}(a_1),\Psi_b(b)\).
These labels suffice for every use of the anchors.

Finally the centralizer threshold defines \(\mathcal B\) in both
\(G_0\) and \(G_\eta\), so every isomorphism between them preserves
this sector. The vector addition just constructed is a group term,
so such an isomorphism is also \(\F\)-linear. Modulo the center it
sends the pure input \(x\) to \((Cx,\beta(x))\) for linear
\(C:X\to X\), \(\beta:X\to Y\). The induced map \(C\) on the
quotient by \(\mathcal B\) is invertible. Pure inputs commute in
\(G_0\). Therefore
\[
 0=\beta(x')Cx-\beta(x)Cx'+\eta(Cx,Cx')
 \qquad(x,x'\in X).
\]
Taking \(\alpha=\beta C^{-1}\) gives \(\partial\alpha=\eta\).
Conversely the shear for \(\partial\alpha=\eta\) is an isomorphism,
with inverse obtained by subtracting \(\alpha(x)\).
This proves~\eqref{hyb:zero-isomorphism}.
\end{proof}

\begin{corollary}[Retained anchors and arbitrary responses]
\label{hyb:responses}
Continue a winning position with these four anchor pairs, and normalize
the mate by a global shear to an old twist
\(\varepsilon:(X^\circ)^2\to Z^\circ\), extended by zero on the
added coordinates, using Proposition~\ref{hyb:cokernel}.
For independent old query vectors \(x_1,\ldots,x_j\in X^\circ\), pebble
\((x_i,0,0)\) in the reference group.  Their responses have quotient
coordinates exactly \(x_i\).  If \(y_i\in Y^\circ\) is the old-output,
old-input part of the response's \(Y\)-coordinate, then
\begin{equation}
 y_i(x_h)-y_h(x_i)=\varepsilon(x_i,x_h).
 \label{hyb:response-equation}
\end{equation}
These queries and their checks require
\(\max\{40,j+4\}\) pairs. Any additional probe is counted separately.
\end{corollary}
\begin{proof}
A global shear fixes the three \(Y\)-anchors and the
center.  It may move the seed's lift, but not its quotient class.
Hence \eqref{hyb:horner}--\eqref{hyb:matrix-name}, including the old
\(X^\circ\)-coordinates and old \(Z^\circ\)-projection, are unchanged.
Their bounded-variable tests force each response's quotient label.
Projecting the commutation of two responses to the old output gives
\eqref{hyb:response-equation}. Central offsets play no role.
If a label or pair relation failed, Spoiler could discard unrelated
queries and test it with the counted variables and a spare pair.
This hypothetical test adds no permanently occupied queries.
No affine form of an announced bijection is assumed.
\end{proof}

\section{Decoding and the reduction}
\label{sec:reduction}

Keep the four anchors from the reconstruction lemma.
After its fixed coordinate identification and a shear, every equivalent
mate has the form $G_{\eps_\beta}$, by
Proposition~\ref{hyb:cokernel}. A shear fixes the center and the
quotient $X$, so all coordinate formulas retain their meaning.
Only the quotient of the seed and the matrix classes of the other
anchors enter those formulas. Their central coordinates are irrelevant.

\subsection{What a frame reads}

For an equation or readout vertex $v$, its frame consists of the $q$
pure group elements $(e_{va},0,0)$, where $e_{va}$ is the scalar unit
in coordinate $a$ of $R_v^q\le X^\circ$. Their answers have the same
$X$-coordinates. Let $y_{va}:X^\circ\to Z^\circ$ be the old-output
component of the answer's matrix coordinate, restricted to $X^\circ$.
It is orthogonal to $C^\circ$ under the matrix pairing. Since all pure
$X$ queries commute in $G_0$, any two retained answers satisfy
\begin{equation}
 y_{va}(e_{wb})-y_{wb}(e_{va})
       =\eps_\beta(e_{va},e_{wb}).
 \label{eq:decoded-commuting}
\end{equation}
The added input and output spaces do not contribute to this equation.
No linearity of Duplicator's announced bijection is assumed.

For an incidence $ij$, define the values read by its two endpoint frames:
\begin{equation}
 \theta_{i\to j}=-\sum_{a,b=1}^q
       \bigl(y_{ia}(e_{jb})\bigr)_{ij;ab,\rad},\qquad
 \theta_{j\to i}=-\sum_{a,b=1}^q
       \bigl(y_{jb}(e_{ia})\bigr)_{ij;ab,\rad}.
 \label{eq:theta}
\end{equation}
The subscript selects the radical coefficient of the indicated
$R_{ij}$ output.

\begin{lemma}[Frame decoding]
\label{lem:decoding}
Every retained equation frame satisfies
$\sum_{j\in T_i}\theta_{i\to j}=\beta_i$.
Every retained readout frame determines a vector $z_j\in V_j$ with
$\theta_{j\to i}=\lambda_{ij}(z_j)$ for all its incidences.
Coexisting endpoint frames have $\theta_{i\to j}=\theta_{j\to i}$.
\end{lemma}
\begin{proof}
Pair $y_{ia}$ with $M_\omega(\tau_i e_{ia})\in C^\circ$ and sum over $a$.
The internal contribution is
\[
 \sum_{a<b}\bigl(y_{ia}(e_{ib})-y_{ib}(e_{ia})\bigr)_{i;ab,\rad}
       =\beta_i,
\]
by~\eqref{eq:decoded-commuting}. Only the pair $(1,2)$ carries a twist.
The incidence contribution is $-\sum_j\theta_{i\to j}$, proving the
equation identity. The tested vectors lie in $L$ because they are radical.

For a readout vertex take any $u=(u_i)_i\in N_j$, pair each $y_{jb}$
with $M_\omega(u e_{jb})$, and sum over $b$. Its internal terms cancel
by~\eqref{eq:decoded-commuting}, since that vertex has no twist.
The incidence terms give $\sum_i u_i\theta_{j\to i}=0$.
Thus $(\theta_{j\to i})_i\in N_j^\perp$, which is exactly the image of
$z\mapsto(\lambda_{ij}(z))_i$. Fix a linear right inverse on this
image to choose the decoded $z_j$. Finally the incidence twist is zero,
so summing~\eqref{eq:decoded-commuting} gives equality of the two values.
\end{proof}

If two equation frames share a scalar readout with nonzero coefficients
at both occurrences, retain both and query that readout's frame. When its coefficient
in equation $i$ is $a_{ij}\ne0$, its underlying value is
$\theta_{i\to j}/a_{ij}$. The lemma forces both values to equal the
decoded readout. This hypothetical continuation uses $3q$ frame pairs
and proves consistency of the existing position. It adds no permanent
frame. The same test forces a zero-coefficient incidence value to vanish.
For each site separately, query its equation frame and its owned
readout frame. Its equation identity then gives zero site shift.

\subsection{From a winning group game to scalar images}

Write $Q=q(r+1)$. Suppose the winning game has at least
$Q+40$ pairs in total, counting the four anchors. For a type $t$ and an allowed full address $a$,
query the $r$ horizontal readout frames of its ring. Let $Q_{t,a}$ be
the set of all sums of their decoded vectors arising in continuations
of the fixed anchoring history that use pure unit queries and removals,
keep the anchors, and end
with only this ring and the anchors retained. Allow every assignment
of frames to pebble slots. These sets are nonempty.


Fix one link and an address permitted on that link. Starting with any
history giving a departure value \(h_t=\sum_i z_i\), replace the
horizontal frame in row \(i\) by its square frame, for each \(i\) in
turn. Place the new frame before releasing the old frame and reuse the
released slots at the next row. When the two coexist, the left horizontal
incidence value of square \(i\) is \(L_{\ell,t}z_i\).
Now all square frames are present. At the vertical readout in row \(i\),
the two adjacent squares have respective coefficients \(+1\) and
\(-1\). The three-frame test therefore makes their two values sum to
zero, including across the last and first rows. The tests can be
performed in separate hypothetical continuations since each checks a
value already fixed by square responses. Replace each square by its
arrival horizontal frame in the same manner. For its decoded value
\(z'_i\), the previously fixed right horizontal incidence value is
\(-L_{\ell,u}z'_i\). Summing the equation identities of the one
simultaneously retained square ring yields
\[
 L_{\ell,t}h_t-L_{\ell,u}h_u=\sum_i\beta_{B(i,\ell,a)}
   =\Delta_\ell(a),\qquad h_u=\sum_i z'_i.
\]
The arriving ring is a continuation from the original anchoring
history, so \(h_u\) is in its defined set. Starting instead from any
arrival ring and reversing the replacements gives the reverse inclusion
of scalar images. At internal wire nodes both forms are the identity,
hence composing the link-image equalities yields~\eqref{hs:images}.
At most $r+1$ frames and four anchors are retained in an actual
transfer. For a consistency or coordinate test, release unrelated
marks and retain the relevant frames or elements together with the
anchors. Since $r\ge3$, $3q+40\le Q+40$, and the formulas of the
reconstruction lemma fit in the same budget. A spare pair remains for
the after-pickup test. This justifies every asserted constraint on a
winning history.

\subsection{The two directions and their overlap}

For an input $w$ to a fixed exponential-time machine, construct its
affine system, and choose
\begin{equation}
 K=r+1,\qquad q=r+100,\qquad Q=q(r+1),\qquad k=Q+40.
 \label{eq:parameters}
\end{equation}
Output the dense alternating tensor of the padded group $G_0$ and $k$.

If the machine accepts, let $H\equiv_k^{II}G_0$. The reconstruction
lemma gives $H\cong G_\eps$ with the prescribed evaluation space.
Proposition~\ref{hyb:cokernel} and a shear replace $\eps$ by
$\eps_\beta$. Keep the four anchors and transport the winning strategy
through these fixed isomorphisms. The $k+1$-pair game has enough space
for the decoder. It forces zero site shifts and the image sets
\eqref{hs:images}. Lemma~\ref{hs:source} gives $\beta\in\im A$.
Hence $\eps_\beta\in\im\partial$, so $H\cong G_0$.
This proves identification against every finite group.

If the machine rejects, take the explicit $\beta^*$ of
Lemma~\ref{hs:source}. It is insoluble and has exact local families
through $K$ blocks. The zero-twist isomorphism criterion in the
reconstruction lemma gives $G_0\not\cong G_{\eps_{\beta^*}}$.
Theorem~\ref{hyb:forward} supplies equivalence for every number of pairs
less than $(q-1)(K+1)$. In particular the pair is Version-II equivalent
through dimension
\begin{equation}
 (q-1)(K+1)-2=Q+96.
 \label{eq:no-range}
\end{equation}
This includes $k$ and leaves room for the version comparisons below.

Lemma~\ref{hyb:local} gives $Z(G)=G'=Z$ for both groups.
The incidence system contains an equation and $q\ge4$, so $Z\ne0$.
Thus each group has class exactly two and exponent $p$.

\paragraph{Versions I and III.}
Set $k_{II}=k$, $k_I=k+6$, and $k_{III}=\lceil k/2\rceil+2$.
The implications between versions are proved in
Lemma~\ref{lem:version-comparisons}. On acceptance, Version-I
equivalence at $k_I$ or Version-III equivalence at $k_{III}$ implies
Version-II equivalence at $k$, so the argument just given forces
isomorphism. On rejection, Version-II equivalence at $Q+96$ implies
Version-I equivalence at $k_I$, since $k_I\le Q+96$.
In turn this implies Version-III equivalence at $k_{III}$, since
$2k_{III}+1\le k_I$ and $k_{III}\ge2$. The same pair thus witnesses
failure of identification at the required dimension in each version.

\paragraph{Equivalence of nonisomorphic pairs.}
Always output $(G_0,G_{\eps_{\beta^*}})$ at the corresponding dimension.
The source shift $\beta^*$ is insoluble on every input, so the two
groups are always nonisomorphic. They are distinguished on acceptance
and equivalent on rejection. Since EXPTIME is closed under complement,
this is an EXPTIME-hardness reduction for equivalence. The upper bound
below completes both claims of Theorem~\ref{thm:main}.

\paragraph{Writing the reduction.}
For a fixed source machine, the compiler gives a polynomial-size list
of circuit types and connection schemas, with their coordinate domains
and bit maps. The vertices and incidences of Section~\ref{hs:section}
are enumerated from types, wires, links, rows and at most two bits per
local label. Domain membership and \eqref{hs:square} determine the
scopes, coefficients and charge. These enumerations never expand the
list of full bit addresses. Fix orders of disjoint summands and
enumerate each finite product of indices lexicographically.

For an explicit size bound, let $s$ be one plus $q$, the numbers of
equations, readouts and incidences, and the sum of the readout
dimensions, all added together. This is polynomial in the source
input length. At a readout $j$ the dimension of $N_j$ is at most
its number of incident equations. Hence $\sum_v\dim R_v=O(s)$,
$\dim X^\circ=O(qs)$ and $\dim Z^\circ=O(q^2s)$, the last bound
including both kinds of output. Write $x=\dim X$ and $z=\dim Z$.
After padding $x,z=O(q^2s)$, $\dim L\le x$, and $\dim Y\le xz$.
The matrix $H$ of Section~\ref{sec:gaussian-bases} therefore has at
most $x$ rows and $xz$ columns. Its entries are computed from the stored
local basis columns by \eqref{hyb:omega}. The local eliminations and
this final elimination take polynomially many field operations on
polynomially many stored entries. The prime $p$ is fixed, so each
field operation can be implemented by a fixed finite table.

For each output tensor, write its two dimensions in unary, each
terminated by a zero. Emit its coefficients in input--input--output
order using \eqref{eq:computed-bracket}, including the diagonal and
both orders of each input pair. Represent a field element by its least
residue $u\in\{0,\ldots,p-1\}$ and emit $1^u0$. There are at most
$(x+xz)^2z$ entries. Identification emits the base description, and
equivalence emits the base followed by the charged description.
Append the positive dimension prescribed for the chosen version, in
binary without leading zeroes or as the unterminated unary word $1^k$.
The dimension is polynomially bounded by \eqref{eq:parameters} and
the displayed version comparisons.

A deterministic multitape transducer can implement the stated loops
with counters and polynomial-length tables. Even scanning an entire
table per access takes polynomial time. To obtain a one-tape machine,
use one track for each simulated tape and mark its head. A constant
number of sweeps of the occupied interval suffices for each step,
and after $t$ steps the interval has length $O(t+|w|+1)$. A final
sweep leaves only the encoded output at the origin. Thus for each
fixed machine, prime, version and dimension encoding the reduction
runs in polynomial time on a deterministic one-tape machine.

\begin{proof}[Proof of Corollary~\ref{cor:unbounded}]
Take the affine system consisting of the single equation $0=1$ and
no readouts. Through $K=0$ blocks its exact local family consists of
the unique empty assignment on the empty set. Given $k\ge1$, put
$q=2k+6$ and take the groups $G_0,G_{\varepsilon_\beta}$ of
\eqref{hyb:group}. Since $\beta=1\notin\im A=\{0\}$,
Proposition~\ref{hyb:cokernel} gives $\varepsilon_\beta\notin\im\partial$.
Thus $G_0\not\cong G_{\varepsilon_\beta}$ by \eqref{hyb:zero-isomorphism}.
Both groups have exponent $p$ by \eqref{hyb:group} and have center
equal to their nonzero derived subgroup by Lemma~\ref{hyb:local}, so
both have class two. Since
$2k+4<q-1=(q-1)(K+1)$, Theorem~\ref{hyb:forward} and
Lemma~\ref{lem:paper-wl-game} give equivalence in Version II at
dimension $2k+3$. Lemma~\ref{lem:version-comparisons} gives
equivalence in Versions I and II at $k$, and in Version III at $k$
for $k\ge2$. When $k=1$, the two groups have equal order, which yields
unary Version-III equivalence by Lemma~\ref{hu:full-gadget}.
\end{proof}

\section{The exponential-time upper bound}
\label{hu:section}

Write $L$ for the input length. A tensor whose input and output spaces
have dimensions $s,t$ represents a group of order $n=p^D$, where
$D=s+t\le L$. Read its two self-delimiting dimensions and its $s^2t$
field entries. Reject an out-of-range entry, a nonzero diagonal entry,
or a failure of skew symmetry. These tests are sufficient, since the
array extends uniquely by bilinearity. For a pair, the count $s^2t$
determines where the second tensor starts. Check the remaining bits
for the prescribed canonical encoding of a positive $k$.
The input tests take polynomial time in $L$, including when a
dimension is zero. For a supplied pair, unequal sums of dimensions give
unequal orders and can be rejected before either table is expanded.
Listing the $n$ coordinate vectors and applying
\eqref{eq:tensor-group} gives the entire Cayley table in
$p^{O(D+1)}$ time for fixed $p$. In each version the group-tuple
histogram has total multiplicity $n^k$. Thus equivalent groups have
the same order, including when just one of them is tensor-described.
In the Version~III graph, each ordered pair $(g,h)$ has its own
four-vertex path $u_0u_1u_2u_3$, with additional edges from
$u_0,u_1,u_3$ to $g,h,gh$, respectively.

\begin{lemma}[A cap on the dimension]
\label{hu:cap}
Every group of order $p^D$ is identified among all finite groups by
each of the three versions at every dimension $k\ge D+10$.
\end{lemma}
\begin{proof}
At order one the claim follows from the multiplicity test. Otherwise
choose generators $a_1,\ldots,a_d$ successively outside the subgroup
already generated, until they generate $G$. At every step the
subgroup order increases by a factor at least $p$, so $d\le D$.
Choose for each $g\in G$ a word $w_g$ in positive powers of these
generators representing $g$, with $w_1$ the empty word. Such words
suffice in a finite group.

Use the ternary relation $M(x,y,z)$ for $xy=z$. There is a sentence
asserting that some $x_1,\ldots,x_d$ give the words $w_g$ pairwise
distinct values and satisfy $w_gw_h=w_{gh}$ for every $g,h\in G$.
In a group of order $n$ the sentence holds exactly when
$g\mapsto w_g(x_1,\ldots,x_d)$ is an isomorphism. It holds in $G$.
Here is an explicit variable bound independent of word length.
Write $W_w(u)$ for evaluation of a word at $u$. Use $M(u,u,u)$ for
the empty word, and use
\[
 W_{w x_i}(u)=\exists v\bigl(W_w(v)\land M(v,x_i,u)\bigr)
\]
with two alternating names for the endpoint and predecessor. An
equality $w_gw_h=w_{gh}$ uses three endpoint names $u,v,z$, their
three separate evaluation formulas, and $M(u,v,z)$. Distinctness
uses two endpoints. One further name suffices inside every evaluation
formula, by rebinding the endpoint in the smaller subformula. Thus
$d+4$ names suffice in the multiplication language.

For the multiplication graph use the vertex predicate $P(x)$ meaning
that $x$ has at least four neighbors. For $n\ge2$ it picks out exactly
the group vertices, of degree $3n$, since auxiliary vertices have
degree at most three. Relativize group quantifiers to $P$. Define
$M(x,y,z)$ by requiring four distinct vertices outside $P$, forming
a path in order, with $x,y,z$ adjacent to its first, second and
fourth vertices, respectively. The components induced outside $P$
are exactly the private four-vertex paths. Only the forward
orientation has its second vertex adjacent to a group element, so
this formula defines the intended ordered product, also for repeated
arguments. Its four auxiliary names and a fifth name for the
counting quantifier in $P$ give at most $d+9$ names altogether.

For completeness, equality of round-$r$ joint colors at two
$k$-tuples preserves the truth of every counting formula of
quantifier depth at most $r$ using their $k$ coordinate names, when
round-zero colors record the atoms. This follows by induction:
Boolean combinations preserve the assertion, and the multiset of
joint replacement vectors, projected to the quantified coordinate,
gives the number of assignments with each preceding color.
Version~I records the multiplication atoms, Version~II records them
as word relations, and the graph records adjacency. Equality of the
restricted histograms supplies tuples of equal round-$r$ color for
each $r$, as the group is nonempty. Take $r$ equal to the depth of
the sentence just constructed. At $k\ge D+10$ it has the same
truth value in the two groups in each version, proving the claim.
\end{proof}

\begin{lemma}[Full histograms on multiplication graphs]
\label{hu:full-gadget}
For $k\ge2$, Version~III equivalence is equivalent to equality, at
each round, of the color histograms on \emph{all} $k$-tuples of the
two multiplication graphs. At $k=1$, Version~III equivalence is
equivalent to equality of group orders.
\end{lemma}
\begin{proof}
Fix $r$ and start from group tuples with matching round-$(r+k)$
colors, which exist when the restricted histograms agree. At each
of $k$ successive replacements, equality of the current colors
gives a bijection of the vertex sets matching the entire preceding
replacement vector, by matching the fibers of every vector. Replace
coordinates $1,\ldots,k$ in order. The bijections may depend on the
preceding choices, and together they give a bijection on all final
vertex tuples preserving their round-$r$ colors.

Conversely, equality of the full histograms gives equal graph orders,
hence equal group orders since the graph has $n+4n^2$ vertices.
Order one is immediate. For $n\ge2$ the round-one color of a tuple
determines which coordinates satisfy $P$. Indeed, keep the coordinate
in question, replace a different coordinate, and count adjacent
replacements using the initial colors. Full histogram equality at
round $r+1$ permits restriction to the tuples all of whose entries
satisfy $P$. Project to round $r$ to recover the required histogram.

In dimension one the graph uses ordinary neighbor refinement. The
five vertex types are the group vertices and the four path positions
$0,1,2,3$. A group vertex has $n$ neighbors of each type $0,1,3$.
The path vertices have neighbor types $(G,1)$, $(G,0,2)$, $(1,3)$,
and $(G,2)$, respectively. Induction on the number of rounds shows
that the color of each type depends only on the order. Conversely,
the cardinality of the restricted histogram is the order.
\end{proof}

\begin{lemma}[Enumerating all comparison groups]
\label{hu:enumeration}
For fixed $p$, in time $p^{O((D+1)^3)}$ one can list group tables of
order $p^D$ containing a representative of every isomorphism class.
\end{lemma}
\begin{proof}
Start with the trivial group. At stage $m$ a table $K$ of order
$p^m$ comes with $m$ distinguished generators. Fix for every $h\in K$
one positive generator word for $h$, using breadth-first search from
the identity with an edge for right multiplication by each generator.
Choose $m$ proposed generator images in $K$ and a carry $c\in K$.
Substitution in the fixed words gives a function $\theta:K\to K$.
On $\{0,\ldots,p-1\}\times K$ form the table
\[
 (a,h)(b,h')=
 \bigl((a+b)\bmod p,\;
      c^{\lfloor(a+b)/p\rfloor}\theta^{\,b}(h)h'\bigr),
\]
where $\theta^{\,b}$ denotes function iteration. Check identity
$(0,1)$, two-sided inverses, and associativity by table lookup.
Check that $(1,1)$ and $(0,s)$ for the distinguished $s\in K$
generate all entries. Retain exactly the tables passing these tests.

For coverage, every nontrivial finite $p$-group of order $p^{m+1}$
has a normal subgroup of index $p$. Indeed the class equation makes
the center's order divisible by $p$. A nonidentity central element
of order $p^a$ has a power of order $p$. Quotient by the resulting
central subgroup and lift a normal subgroup of index $p$ by
induction, with the group of order $p$ as base case. Given this
normal subgroup $K$, let $x$ lift a generator of the cyclic
quotient. Each group element has a unique expression $x^a h$ with
$0\le a<p$ and $h\in K$. Take the actual generator images under
$h\mapsto x^{-1}hx$ and take $c=x^p$. Then
\[
 x^ahx^bh'=x^{(a+b)\bmod p}
    c^{\lfloor(a+b)/p\rfloor}\theta^{\,b}(h)h'.
\]
Induction and transport through a tail isomorphism show that this
choice occurs in the enumeration and passes all tests.

At stage $m$ there are $p^{m(m+1)}$ choices per preceding table.
The number of complete choice sequences is at most
$p^{\sum_{m=0}^{D-1}m(m+1)}\le p^{D^3}$, and the number of prefixes
is at most $D+1$ times that bound. Each table computation, test,
and word search takes polynomial time in $p^{D+1}$. Repetitions do
not affect coverage or the bound.
\end{proof}

\begin{proposition}
\label{hu:membership}
For each version and each encoding of $k$, tensor-input identification
among all finite groups and tensor-input equivalence belong to
$\mathrm{EXPTIME}$.
\end{proposition}
\begin{proof}
After input validation, compare $k$ with $D+10$ before allocating any
tuple table. For instance, scan the parameter with a running value
saturated at $D+10$. Isomorphic groups have matching colors in
every round by transporting replacement choices along the isomorphism.
Thus by Lemma~\ref{hu:cap}, above or at the cap identification is
automatic and equivalence is precisely isomorphism. For a pair first
reject unequal group orders. At equal order an isomorphism test costs
$n^{D+O(1)}$: choose at most $D$ generators of the first group by
successive subgroup searches, fix a representative generator word
for each element, and try every assignment of generator images in
the second group. Test whether substitution yields a bijection
respecting the entire multiplication table. Every isomorphism is
tested and each test is polynomial in $n$.

Below the cap the two structures are refined together with shared
color numbers and with replacements in each tuple's own structure.
For Version~II, to compare the initial types of tuples $a,b$,
generate the subgroup of the product generated by the pairs
$(a_i,b_i)$, by a search from the identity. Its projections onto
the two marked subgroups are surjective. They are both injective
exactly when the correspondence $a_i\mapsto b_i$ extends to an
isomorphism of these subgroups, equivalently when the word types
agree. This costs polynomial time in $k$ and $n$ for each comparison.
Within one structure the same test applies. Version~I only needs
the equality and multiplication tables of tuple entries. For
Version~III use Lemma~\ref{hu:full-gadget}.

The total number of tuples is at most $Q=2(n+4n^2)^k$.
Each color partition refines its predecessor and a round without a
split fixes all subsequent partitions, since the new labels then
only rename the old classes. Thus at most $Q$ strict refinements
are possible. A round can be carried out in polynomial time in
$Q,n,k$ by sorting the replacement vectors and giving equal
signatures the same integer label. The stable histogram answers
equivalence since each color retains its parent. For $k<D+10$,
this computation takes $p^{O((D+1)^2)}$ time.

For identification below the cap, apply the equivalence test and
the isomorphism test to each table from Lemma~\ref{hu:enumeration}.
Reject exactly when some equivalent table is nonisomorphic to the
input. The order test proves that this quantifies over all possible
finite comparison groups. For a supplied pair only its own table
comparison is needed. The overall running time is
$L^{O(1)}+p^{O((D+1)^3)}$ for fixed $p$. A polynomial overhead for
single-tape simulation preserves this bound. Since $D\le L$, both
decision problems have deterministic exponential-time algorithms.
\end{proof}

Together with the polynomial-time reductions of Section~\ref{sec:reduction},
Proposition~\ref{hu:membership} proves Theorem~\ref{thm:main}.

\appendix

\section{Compiling a bounded machine into address schemas}
\label{app:compiler}
An addressed circuit with $r$ bits has one gate $(g,a)$ for each
gate type $g$ and $a\in\{0,1\}^r$. A connection schema specifies
two types, a target input port, a nonempty product of subsets of
$\{0,1\}$ as its source domain, and one bit permutation per coordinate.
It joins each address in the domain to its image. A scalar gate is
the OR of its seeds and inputs. An AND gate is true when each of its two
ports has a true incoming gate. Seeds and output tests specify a scalar
type and a product domain. The circuit accepts if a tested gate is true.

\begin{lemma}
\label{lem:machine-schemas}
For each deterministic one-tape exponential-time decider $M$ there is a
polynomial-time transducer that, on input $w$, writes an acyclic
addressed circuit accepting exactly when $M$ accepts $w$. The number
$r$ of address bits is polynomial in $|w|$ and may be taken at least
$32$. The number of gate types is constant for fixed $M$, the number of
connection schemas is $O(r^2)$, and the total description length is
$O(r^3)$.
\end{lemma}
\begin{proof}
Fix integers $c,e$ such that $M$ halts on length-$m$ inputs within
$2^{c(m+1)^e}$ steps. Increasing $c$ if necessary, put
\[
 d=c(m+1)^e+(m+1)+2,\qquad T=2^{d-2},\qquad N=2^d,\qquad r=2d\ge32 .
\]
The bound $T$ covers both the running time and $m+1$.
Use a circle of $N$ tape cells and rows $t=0,\ldots,N-1$. A letter
consists of the tape symbol and either no head or the control of the
head. The alphabet $\Gamma$ is fixed. In a legal row the center letter
at the next time is a function $f$ of the three neighboring letters.
The old head writes the prescribed symbol, and the move determines the
only new head. A halted head stays in place and preserves its symbol.
Extend $f$ arbitrarily to triples incompatible with a legal row.
Initialize the circle with the input in cells $0,\ldots,m-1$,
blanks elsewhere and the initial head at zero. Until the halt time
the head is in $[-T,T]$, which embeds injectively in the circle and
contains the written input. The initial contents agree on that interval.
Induction on the number of executed moves gives the same symbol read,
transition, write and next head position as on the integer tape.
After the halt the entire legal row is fixed. Thus the last row
contains an accepting head exactly when $M$ accepts.

Use $d$ bits for $t$ and $d$ bits for the spatial index $x\bmod N$.
For each letter $a$ use one scalar type $S_a$. For each
$(a,b,c')\in\Gamma^3$ use two AND types $P_{abc'}$ and $U_{abc'}$.
At $(t+1,x)$ the inputs to $P_{abc'}$ are
$S_a(t,x-1)$ and $S_b(t,x)$ on different ports.
At the same address the inputs to $U_{abc'}$ are $P_{abc'}(t+1,x)$
and $S_{c'}(t,x+1)$ on different ports.
For every address join $U_{abc'}$ to $S_{f(a,b,c')}$ there.
The same-address wire from $P$ to $U$ is also present in row zero.
The $P$ gates in that row have no inputs from an earlier row.
Seed precisely the scalar letter prescribed at time zero.
Give $P,U,S$ in row $t$ ranks $3t,3t+1,3t+2$, respectively.
Every wire increases rank. Induction in this order shows that exactly
one scalar letter is true at each address, namely the corresponding
letter in the cellular tableau. At positive times exactly the triple
of letters in the preceding neighborhood reaches its $U$ gate.
Tests inspect all scalar types carrying an accepting head at time
$N-1$ and arbitrary $x$.

To encode a nonwrapping increment on $d$ bits, take one case for each
possible least zero bit $j$. Fix all bits below $j$ to one and bit $j$
to zero, and flip bits $0,\ldots,j$. Every case is a product domain.
For a cyclic spatial increment also allow the all-ones string and
flip all bits. Decrement uses the reversed bit conditions.
The two halves of the address are disjoint, so a simultaneous
time and spatial shift uses the products of their cases. For each
triple this gives $d(d+1)$ schemas for the left input, $d$ for the
center, $d(d+1)$ for the right, and two for the same-address wires.
Omitting the time wrap prevents any input from the final row to the
initial row.

For the seeds, write individual requests for positions $0,\ldots,m$
in the initial row. This handles an empty input and the head at zero.
The remaining blank positions form the interval $[m+1,N-1]$.
Use the singleton $m+1$ and, for each zero bit of $m+1$, the product
domain where the higher bits agree with $m+1$, that bit is one,
and lower bits are free. The highest differing bit proves that these
products cover the interval exactly. The compiler can enumerate all
gates, triples, pairs of carry positions, and address bits with fixed
nested loops. It computes the polynomial integer $d$ and the $d$-bit
word for $m+1$. It never enumerates $N$ positions or $N$ times.
Since $m<d$, there are $O(r)$ seed and test requests and $O(r^2)$
connections, each with $O(r)$ bits of coordinate data.
Unary delimited record counts also have polynomial length.
All loops and all work strings have polynomial bounds, so the
description is computable in polynomial time. A one-tape transducer
can perform the same operations by scanning its polynomially bounded
work area for each simulated operation.
\end{proof}

The later cyclic layout needs $r\ge3$, equivalently $r=n+3$ with
$n\ge0$. Here $n=2d-3$. No addresses need to be appended to the
compiled circuit. In the identification reduction we apply the lemma
to a machine for the source language. For the equivalence reduction we
interchange its two halting answers before compilation, since the
constructed pair is equivalent precisely on rejection.

\section{Refinement and the three versions}
\label{app:games}

Call two tuples \emph{matched} if their colors agree at every round for
the joint update of Section~\ref{sec:statement}. Write $a[i/x]$ for a
coordinate replacement. For matched tuples $a,b$ there is a single
bijection $f$ between the underlying sets such that
$a[i/x]$ matches $b[i/f(x)]$ for every $i,x$. Indeed at any given round
the replacement vectors have matching multiplicities. Pair within their
fibers. The sets of admissible bijections for successive rounds are
nested nonempty finite sets. Their intersection supplies $f$.
Similarly, equality of histograms at all rounds supplies a matching
tuple for any prescribed tuple in the histogram. Conversely one pair
of matched tuples gives a bijection between the full tuple sets
preserving stable colors: replace their coordinates in order using the
announced bijection conditional on the preceding choices. Each such
conditional map is bijective, so the induced map on sequences is
bijective. The same argument applies on a subset when each successive
bijection carries that subset onto its counterpart.

\begin{lemma}[The Version-II game]\label{lem:paper-wl-game}
For $d\ge1$, Version-II equivalence is equivalent to a winning
Duplicator strategy, from the empty position, with $d+1$ pebble pairs.
On each move Spoiler picks up one pair (possibly unused), the remaining
marks must induce an isomorphism of their generated subgroups, Duplicator
announces a bijection, and Spoiler places the pair at an element and
its image or chooses the image and uses its preimage.
A winning position with $m+1$ pairs preserves every counting formula
on its marked parameters that uses at most $m$ variable names, with
relations between group words as atomic formulas.
\end{lemma}
\begin{proof}
Assume the histograms agree. The all-identity tuples match, since
their initial word type forces each entry of a matching tuple to be the
identity. Begin with $d+1$ virtual identity pairs. Maintain that deleting
any one pair gives matching $d$-tuples in corresponding order.
When pair $i$ is removed its complement matches, so the subgroup test
passes. Apply the replacement bijection of that complement. After
placement, omitting $i$ still gives the old complement. Omitting any
other pair gives a coordinate replacement of the old complement up to
a common reordering. A common coordinate permutation preserves matching,
by induction on rounds. Ignoring initially unused virtual marks gives a
strategy on the empty board.

Conversely, take any winning position containing $d$ designated pairs.
Their $d$-tuples have equal colors at every round. At round zero lift
the remaining pair and apply the subgroup test. For the induction step,
lift that same pair. Its announced bijection matches all replacement
vectors: after placing $x$, designate the new pair in place of the old
$i$th pair and use the induction hypothesis, independently for each
$i$. Starting from the empty board, fill $d$ designated pairs in
order. The resulting conditional bijections match the multiplicities
of all tuples. For the last assertion, reserve the extra pair for
testing atoms after pickup. At a counting quantifier move the pair of
the rebound variable. The announced bijection pairs the satisfying
values by induction on the subformula. Negation and conjunction
preserve the assertion. Any marks unrelated to the free variables
can be ignored by maintaining them as virtual marks. Unused variable
names can then be assigned with free pairs.
\end{proof}

\begin{lemma}[Changes of version]\label{lem:version-comparisons}
For finite groups, the following implications hold:
\begin{align*}
 G\equiv^{II}_d H&\ \Longrightarrow\ G\equiv^{I}_k H
     &&(d\ge k\ge1),\\
 G\equiv^{I}_d H&\ \Longrightarrow\ G\equiv^{II}_k H
     &&(d\ge k+2,\ k\ge1),\\
 G\equiv^{I}_d H&\ \Longrightarrow\ G\equiv^{III}_k H
     &&(d\ge2k+1,\ k\ge2),\\
 G\equiv^{III}_d H&\ \Longrightarrow\ G\equiv^{II}_k H
     &&(d\ge\lceil k/2\rceil+2,\ k\ge1).
\end{align*}
Version-II equivalence is also monotone downwards in dimension.
\end{lemma}
\begin{proof}
For the first two implications designate $k$ entries of matching
$d$-tuples. In the first case their Version-I initial colors agree,
because equality and multiplication are word tests. In the second
case all word tests on these entries can be performed with two spare
variable names $s,t$ in the multiplication relation. For a signed word
$L$, test $s=L$ as follows. For the empty word test $s^2=s$. For
$L=x_iL'$ quantify $t$ with $x_it=s$ and recursively test $t=L'$,
reusing $s$ as the spare name. For $L=x_i^{-1}L'$ use $x_is=t$
instead. To test $L=1$, first quantify $s$ with $s^2=s$. Matching
Version-I tuples satisfy the same such formulas: the atomic tests are
in the initial color and each quantified witness is transported by the
replacement bijection. In either case, induction on the output round
now gives equal colors on the designated entries, since one announced
bijection handles every designated replacement. Obtain their full
histograms by filling those entries in order. The same argument with
Version II on both sides proves its monotonicity.

For the third implication write $P_j(x,y)$ for the private vertex in
position $j=0,1,2,3$ of the path for $(x,y)$. Represent each graph
vertex by two distinct group slots and a tag. The group tag reads
only the first slot, and tag $j$ reads $P_j(x,y)$ from both. Reserve
one additional group slot. Matched Version-I tuples with identical
tags induce equal initial graph colors: attachments to a group vertex
$z$ test respectively $z=x$, $z=y$, or $xy=z$. Private edges test
equality of the ordered pair and consecutiveness of the tags.
We show by induction that the induced graph tuples have equal colors
at every round. Choose the matching bijection $f$ for the spare slot.
After placing $x,f(x)$ there choose the matching bijection $f_x$ of
the resulting tuples. Announce the graph bijection
\[
 x\longmapsto f(x),\qquad
 P_j(x,y)\longmapsto P_j\bigl(f(x),f_x(y)\bigr).
\]
For replacement at graph position $i$, overwrite its old first slot
with $y$ using $f_x$. Use the former spare and that slot as the new
pair, and the old second slot as the new spare. For a group vertex the
second slot can hold any matched value. All other representations
remain in place. The induction hypothesis gives the required previous
color in every component of the replacement vector. To start the
histogram, take matching group tuples of dimension $d$, choose group
tags throughout and fill the $k$ first slots in order. This matches
all group-vertex tuples, as required for Version III.

For the fourth implication equality of the restricted histograms gives
$|G|^d=|H|^d$. The common order one case is immediate, so suppose it
is at least two. The five sorts in the graph are visible to matching
$d$-tuples. Group vertices have degree $3|G|>3$, $P_1$ has degree
three and the three other sorts have degree two. Among those three,
$P_2$ has no group neighbor, while $P_0$ has both a group neighbor and
a degree-three neighbor and $P_3$ has only the former. Degree and
existence of a neighbor of any recognized sort are preserved by stable
matching: keep the tested vertex and use a different slot to count or
choose neighbors under its matching bijection.

From $P_1(x,y)$ one can read the three coordinates $x,y,xy$ by
following respectively the sort sequences $(0,\mathrm{group})$,
$(\mathrm{group})$, and $(2,3,\mathrm{group})$. The group vertex $x$
itself may serve as a degenerate representation with all three
coordinates equal to $x$. Equality between two such implicit
coordinates of marked positions is preserved by matching tuples.
Indeed it is a composition of the first typed path with the reverse of
the other, with the zero-length group branch for a degenerate
representation. Any such binary relation built by composition,
intersection, union and reversal from equality, adjacency and sort
tests is preserved using three names. For composition retain the two
endpoints, choose an intermediate vertex with the third name, and
apply induction to the two subrelations. Names no longer free in a
subrelation may be rebound there. Here $d\ge3$.

Pack the $k$ group entries in consecutive pairs of implicit inputs
of $\lceil k/2\rceil$ graph positions, leaving the last second input
unused if $k$ is odd. The two remaining graph names suffice to test
any word in the packed entries. In detail, an output coordinate $s$
represents the empty word precisely when it is the identity. Identity
can itself be tested by equality to an input of a $P_1$ whose three
coordinates coincide, since $a^2=a$ forces $a=1$. This test is a
binary-relation composition and uses only three names locally. To
test $s=x_i L'$ introduce a helper $P_1(x_i,t)$, require its first
input to equal the indicated packed coordinate, its product to equal
$s$, and recursively test its second input $t=L'$. To test
$s=x_i^{-1}L'$ use the helper with second input $s$ and recursively
test its product against $L'$. After either step the former output
name is free for recursion. To start, bind one of the two names to
the identity group vertex as output. This is a formula in the $d$
graph names: the three-name tests of implicit coordinate equality are
separate subformulas and may shadow names not free in that test.
Induction on the word proves the stated semantics. Stable matching
therefore gives the same Version-II initial color to the packed group
tuples.

For a refinement step use the first spare graph position to choose a
new group element $x$. The announced graph bijection restricts to a
bijection $f:G\to H$ because its images preserve sorts after
placement. For a replacement at group coordinate $i$ let $u$ be the
other implicit input of its packed position. Use the second spare
position to mark $P_1(x,u)$ or $P_1(u,x)$ in the appropriate order.
On the other graph its mate must be the corresponding unique $P_1$
with the matched inputs, by the coordinate-equality tests against the
retained positions. Use this new position in the packing and release
the old position. The resulting packed tuples retain every other coordinate. Their
values at the designated coordinate become $x$ and $f(x)$. The bijection $f$ was chosen before $i$, so
the same $f$ matches the whole replacement vector. Induction on rounds
gives stable Version-II matching of all represented tuples. Finally
choose any matching pair of group-only $d$-tuples from the restricted
histograms and initially use their entries as degenerate
representations. One resulting pair of matched $k$-tuples gives equal
Version-II histograms by sequential replacement, as above.
\end{proof}

\begingroup
\small
\begin{thebibliography}{9}
\setlength{\itemsep}{3pt}
\setlength{\parskip}{0pt}

\bibitem{BS2020}
Jendrik Brachter and Pascal Schweitzer.
\emph{On the Weisfeiler--Leman Dimension of Finite Groups}.
Proceedings of LICS 2020, pp.\ 287--300. Full version,
arXiv:2003.13745v2, 2021.
\url{https://arxiv.org/abs/2003.13745}.

\bibitem{OpenAI2026}
OpenAI.
\emph{The complexity of identifying a graph by Weisfeiler--Leman refinement}.
Manuscript, 25 September 2026.
\href{https://github.com/openai/math/tree/fd4aeeb2ee4fc729c18d98444fed42fd0529eeeb/preprints/The-complexity-of-identifying-a-graph-by-Weisfeiler-Leman-refinement-September-25-2026}{Official manuscript and source}.

\end{thebibliography}
\endgroup
\end{document}
